% ch09_conclusion.tex
% Chapter 9: Conclusion, Limitations, and Future Work
% Octopus-Inspired Soft Robotic Drone Navigation Using Bio-Mimetic SLAM

\selectlanguage{hebrew}

\section{סיכום, מגבלות המחקר והצעות להמשך}
\label{sec:ch09_conclusion}

\subsection{סיכום התרומות המרכזיות}
\label{sec:ch09_contributions}

מאמר זה הציג מערכת ניווט ביומימטית מקיפה לרחפן רובוטי רך תת-ימי, בהשראת מערכת העצבים המבוזרת של התמנון. המערכת המוצעת משלבת ארבע תרומות מרכזיות המתכתבות עם האתגרים הייחודיים של ניווט רובוטים רכים בסביבות תת-ימיות מוגבלות.

התרומה הראשונה היא מודל קינמטי ודינמי מקיף לרחפן רך רב-זרועי, המשלב קינמטיקת PCC (Piecewise Constant Curvature) עם דינמיקה מבוססת אלמנטים סופיים וכוחות הידרודינמיים. המודל מאפשר חישוב קינמטי יעיל בזמן אמת עם דיוק של 5.2\% בהשוואה למודל אלמנטים סופיים מלא, תוך שמירה על קצב עדכון של 200 Hz על פלטפורמת ARM משובצת. המודל ההידרודינמי כולל כוחות גרר, מסה נוספת וכוחות ציפה, המאפשרים הדמיה ריאליסטית של תנועת הרחפן בסביבה תת-ימית.

התרומה השנייה היא ארכיטקטורת SLAM מבוזרת בהשראת מערכת העצבים של התמנון, המבוססת על גרף פקטורים מבוזר עם אילוצי קונצנזוס מסוג ADMM. הארכיטקטורה מפרקת את בעיית ה-SLAM הגלובלית לתתי-בעיות עצמאיות לכל זרוע, תוך שמירה על עקביות גלובלית באמצעות תקשורת מבוזרת. המורכבות החישובית של SLAM מבוזר היא $O(N_{\text{arms}})$ לעומת $O(N_{\text{arms}}^2)$ ל-SLAM ריכוזי, מה שמאפשר מדרגיות לרחפנים בעלי 8-12 זרועות.

התרומה השלישית היא שיטת מיזוג חיישנים רב-מודאלית המשלבת סונאר, ראייה סטריאוסקופית, חישת מישוש ו-IMU, תוך שימוש בפילטר חלקיקים מותאם עם שקלול מבוסס אי-ודאות. השיטה מאפשרת למערכת להמשיך בניווט מדויק גם כאשר חלק מהחיישנים אינם פעילים, על ידי העברת משקל דינמית בין מודאליות החישה. מיפוי הסביבה ההיברידי משלב מיפוי סונאר למרחקים גדולים (עד 15 מ') עם מיפוי מישוש ברזולוציה גבוהה (4 ס"מ) למגע קרוב.

התרומה הרביעית היא שיטת תכנון תנועה אדפטיבית המבוססת על למידת חיזוק מבוזרת מרובת-סוכנים (MAPPO), שבה כל זרוע פועלת כסוכן עצמאי תחת מבקר מרכזי. פונקציית התגמול מבוססת על רווח מידע SLAM, קירוב ליעד, חיסכון באנרגיה והימנעות מהתנגשויות. השיטה משיגה שיעור הצלחה של 86\% ואורך נתיב ממוצע של 10.3 מ', שיפור משמעותי לעומת שיטות ייחוס.

המערכת המשולבת הוערכה בשלושה תרחישי סימולציה מייצגים: מערה תת-ימית, שדה אצות ומבנה תעשייתי. התוצאות מראות ATE ממוצע של 5.5 ס"מ, RPE ממוצע של 1.7 ס"מ/מ', ושיעור הצלחה ממוצע של 93\% : שיפור של 36\% ב-ATE לעומת SLAM ריכוזי ו-52\% לעומת ORB-SLAM3 מותאם לסביבה תת-ימית. זמן החישוב הממוצע הוא 18 ms לפריים, המאפשר פעולה בזמן אמת על פלטפורמת ARM Cortex-A72 עם צריכת הספק כוללת של 5.8 W.

הטבלה הבאה מסכמת את התרומות המרכזיות של המחקר וממפה אותן לשאלות המחקר המרכזיות:

\begin{table}[htbp]
\centering
\caption{סיכום התרומות המרכזיות של המחקר ומיפוי לשאלות המחקר.}
\label{tab:ch09_contributions}
\adjustbox{max width=\columnwidth}{%
\begin{tabular}{p{1.5cm}p{4cm}p{4cm}p{3cm}}
\toprule
\textbf{תרומה} & \textbf{תיאור} & \textbf{שאלת מחקר} & \textbf{תוצאה מרכזית} \\
\midrule
1 & מודל קינמטי ודינמי לרחפן רך & כיצד למודל רובוט רך רב-זרועי? & דיוק 5.2\%, 200 Hz \\
2 & SLAM מבוזר בהשראת תמנון & כיצד לבצע SLAM מבוזר? & ATE 5.3 ס"מ, $O(N)$ \\
3 & מיזוג חיישנים רב-מודאלי & כיצד למזג חיישנים הטרוגניים? & ATE 6.3 ס"מ, 94\% הצלחה \\
4 & תכנון תנועה אדפטיבי ב-RL & כיצד לתכנן תנועה אדפטיבית? & 86\% הצלחה, 10.3 מ' נתיב \\
\bottomrule
\end{tabular}%
}
\end{table}

\subsection{מגבלות המחקר : מגבלות חומרה, סביבת ניסוי, והכללה}
\label{sec:ch09_limitations}

למרות התוצאות המבטיחות, למחקר מספר מגבלות משמעותיות שיש להכיר בהן. ראשית, כל הניסויים בוצעו בסימולציה בלבד. למרות שהסימולטור מבוסס על מודלים פיזיקליים מדויקים (SOFA לדינמיקת גוף רך, UUV Simulator להידרודינמיקה), הפער בין סימולציה למציאות (Sim-to-Real Gap) עלול להיות משמעותי עבור רובוטים רכים. אי-ליניאריות חומרית, היסטרזיס, שינוי תכונות החומר עם הזמן והטמפרטורה, ואפקטים של חיכוך לא-מודל : כל אלה עלולים לגרום לסטייה משמעותית בין הביצועים בסימולציה לבין הביצועים במציאות.

שנית, מגבלות חומרה משמעותיות. מערך 96 חיישני המישוש הקיבוליים דורש 12-bit ADC עם 96 ערוצים, מה שמצריך FPGA או microcontroller ייעודי. צריכת ההספק הכוללת של מערכת החישה (480 mW) פלוס מעבד ARM (3.8 W) מגבילה את זמן הפעולה של רחפן המופעל על ידי סוללת 50 Wh לכ-8.6 שעות בתנאי ניסוי סטנדרטיים. במצבים של זרמים חזקים או תמרונים אינטנסיביים, צריכת ההספק עלולה לעלות משמעותית.

שלישית, מגבלות סביבת הניסוי. שלושת התרחישים שנבחנו מייצגים סביבות טיפוסיות אך אינם מכסים את כל מגוון הסביבות התת-ימיות האפשריות. סביבות עם עכירות קיצונית (ראות < 0.1 מ'), זרמים חזקים (> 2 מ'/ש'), טמפרטורות קיצוניות (< 2°C או > 35°C), או מליחות משתנה לא נבחנו. כמו כן, כל הניסויים בוצעו בסביבה סטטית ללא תנועה של עצמים בסביבה.

רביעית, מגבלות הכללה. הארכיטקטורה המבוזרת תוכננה במיוחד עבור רחפן בעל 6 זרועות בתצורה סימטרית. הכללה לרחפנים בעלי מספר שונה של זרועות או בתצורה לא-סימטרית דורשת התאמה של פרמטרי האלגוריתם. כמו כן, המודל הקינמטי מבוסס על הנחת PCC המתאימה לזרועות סיליקון גליליות, אך עשויה שלא להתאים לזרועות בעלות גאומטריה מורכבת יותר.

חמישית, מגבלות זמן אמת. למרות שהמערכת פועלת בקצב 50 Hz על פלטפורמת ARM, עומס חישובי גבוה במצבים מסוימים (למשל, סגירת לולאה עם מספר רב של תכונות) עלול לגרום להחמצת דדליין. מנגנון עדיפויות והשלכת פריימים נדרש כדי להבטיח פעולה יציבה.

\subsection{כיווני מחקר עתידיים : רובוט רך אוטונומי לחלוטין, למידה אנד-טו-אנד, מערכות רב-רחפנים}
\label{sec:ch09_future}

שלושה כיווני מחקר עתידיים עיקריים מזוהים, המבוססים על המגבלות שזוהו ועל המגמות העדכניות בספרות.

הכיוון הראשון הוא פיתוח רובוט רך אוטונומי לחלוטין, שבו כל רכיבי המערכת : חישה, עיבוד, הסקה ובקרה : משולבים בגוף הרך עצמו ללא רכיבים קשיחים. גישה זו דורשת פיתוח חיישנים ומעבדים גמישים הניתנים להטמעה בחומר הסיליקון. התקדמויות אחרונות בתחום האלקטרוניקה הגמישה (Shih et al., 2020) מראות כי ניתן להטמיע חיישנים קיבוליים, טרנזיסטורים גמישים ואפילו מעבדים פשוטים ישירות בחומר הרך. אתגר משמעותי הוא פיתוח סוללה גמישה או מקור אנרגיה חלופי (כגון תא דלק ביולוגי) המסוגל לספק את צריכת ההספק הנדרשת.

הכיוון השני הוא מעבר ללמידה אנד-טו-אנד (End-to-End Learning), שבו רשת עצבית אחת לומדת ישירות מחיישנים לפעולות, תוך עקיפת שלבי הביניים של SLAM ותכנון תנועה. גישה זו עשויה לשפר את הביצועים בסביבות מורכבות על ידי למידה ישירה של אסטרטגיות ניווט אופטימליות מנתוני אימון. עם זאת, גישה אנד-טו-אנד דורשת כמויות גדולות של נתוני אימון (מיליוני פריימים) ומערכת תגמול מתוחכמת. כמו כן, היעדר ייצוג מפורש של מפת הסביבה מקשה על הבנת החלטות המערכת (explainability) ועל טיפול במצבי כשל.

הכיוון השלישי הוא הרחבה למערכות רב-רחפנים, שבהן מספר רחפנים רכים משתפים פעולה במיפוי וניווט בסביבה משותפת. ארכיטקטורת ה-SLAM המבוזרת שהוצגה במאמר זה מהווה בסיס טבעי להרחבה כזו, כאשר כל רחפן מהווה "גנגליון-על" בארכיטקטורה מבוזרת רחבה יותר. אתגרים משמעותיים כוללים תקשורת בין-רחפנים בסביבה תת-ימית (מוגבלת לטווח קצר ורוחב פס נמוך), חלוקת משימות דינמית, וטיפול בקונפליקטים בין מפות חלקיות של רחפנים שונים.

כיווני מחקר נוספים כוללים: (1) שילוב של מודל חומר אדפטיבי המתעדכן בזמן אמת על בסיס מדידות חיישנים, לצמצום פער המודל; (2) פיתוח מנגנון חקר מובנה (Structured Exploration) המשלב סריקה שיטתית עם חקר מבוסס אי-ודאות; (3) שילוב של חיזוי מזג אוויר וזרמים ימיים בתכנון התנועה; (4) פיתוח ממשק אדם-מכונה המאפשר למפעיל אנושי לקבל החלטות ברמה גבוהה תוך השארת השליטה המקומית למערכת המבוזרת.

\subsection{סיכום ומסקנות}
\label{sec:ch09_summary}

לסיכום, המחקר הנוכחי מהווה צעד משמעותי לקראת מימוש חזון הרובוט הרך האוטונומי המסוגל לנווט בסביבות תת-ימיות מורכבות. השראה ביולוגית ממערכת העצבים של התמנון, בשילוב עם כלים חישוביים מודרניים מבוססי גרפי פקטורים ולמידת חיזוק, מאפשרת התמודדות עם אתגרי הניווט הייחודיים של רובוטים רכים בסביבות מוגבלות.

המערכת המוצעת מדגימה כי ארכיטקטורה מבוזרת בהשראת ביולוגיה יכולה להתעלות על גישות ריכוזיות מסורתיות במונחים של דיוק, עמידות לכשל, מדרגיות ויעילות חישובית. השיפור הממוצע של 36\% ב-ATE לעומת SLAM ריכוזי, והפחתת העומס החישובי ב-60\%, מדגימים את הפוטנציאל של גישה זו.

האתגר המרכזי להמשך הוא המעבר מסימולציה למציאות. פיתוח אב-טיפוס פיזי של הרחפן הרך, עם כל מערך החיישנים המבוזר ומערכת העיבוד המוטבעת, יאפשר אימות של התוצאות בסביבה תת-ימית אמיתית. אתגר זה דורש שיתוף פעולה בין-תחומי בין מהנדסי רובוטיקה רכה, מומחי חיישנים, חוקרי למידת מכונה וביולוגים ימיים.

\medskip
\noindent\textbf{תודות:} המחקר נתמך על ידי הקרן הלאומית למדע (ISF) ומשרד המדע והטכנולוגיה.

\appendix
\section{רשימת סמלים ומשתנים}
\label{sec:ch09_appendix}

\begin{table}[htbp]
\centering
\caption{רשימת סמלים ומשתנים בשימוש במאמר.}
\label{tab:ch09_symbols}
\adjustbox{max width=\columnwidth}{%
\begin{tabular}{p{2cm}p{5cm}p{3cm}p{2cm}}
\toprule
\textbf{סמל} & \textbf{הגדרה} & \textbf{יחידות} & \textbf{ערך טיפוסי} \\
\midrule
$\mathbf{x}_t$ & וקטור מצב הרחפן בזמן $t$ & משתנה & : \\
$\mathbf{m}$ & מפת הסביבה & משתנה & : \\
$\mathbf{z}_t$ & וקטור תצפית בזמן $t$ & משתנה & : \\
$\mathbf{u}_t$ & וקטור בקרה בזמן $t$ & משתנה & : \\
$\kappa_i$ & עקמומיות מקטע $i$ & m$^{-1}$ & 0-5 \\
$\phi_i$ & זווית מישור כיפוף מקטע $i$ & rad & 0-$2\pi$ \\
$l_i$ & אורך קשת מקטע $i$ & m & 0.1-0.4 \\
$\mathbf{T}_i$ & טרנספורמציה הומוגנית של מקטע $i$ & מטריצה $4\times4$ & : \\
$\rho$ & צפיפות הנוזל & kg/m$^3$ & 1025 \\
$C_d$ & מקדם גרר & חסר ממד & 0.8-1.2 \\
$\mathbf{F}_{\text{drag}}$ & כוח גרר הידרודינמי & N & 0-0.5 \\
$C_{ij}$ & קיבול חיישן מישוש בתא $ij$ & F & 1-10 pF \\
$\varepsilon_r$ & מקדם דיאלקטרי יחסי & חסר ממד & 2.5-3.5 \\
$\Sigma_t^{\text{SLAM}}$ & מטריצת קווריאנס SLAM & משתנה & : \\
$\text{ATE}$ & שגיאת מסלול מוחלטת & m & 0.04-0.12 \\
$\text{RPE}$ & שגיאת תנוחה יחסית & m/m & 0.01-0.05 \\
$\Delta H$ & הפחתת אנטרופיית מפה & bits & 0.3-0.6 \\
$N_{\text{arms}}$ & מספר זרועות הרחפן & חסר ממד & 6 \\
$\rho_{\text{ADMM}}$ & פרמטר עונשין ADMM & חסר ממד & 1.0 \\
$\gamma$ & פקטור היוון ב-RL & חסר ממד & 0.99 \\
$\epsilon$ & פרמטר גזירה PPO & חסר ממד & 0.2 \\
$\alpha_1,\alpha_2,\alpha_3,\alpha_4$ & משקלי פונקציית תגמול & חסר ממד & 1.0, 0.5, 0.1, 2.0 \\
\bottomrule
\end{tabular}%
}
\end{table}

\cite{Thrun2005, Dellaert2017, Sumanaweera2026, Laschi2016, Cadena2016, Rus2015, Schulman2017, Lowe2017, Mazzolai2020}
